% ═══════════════════════════════════════════════════════════════════════════
\section{Istruzioni per l'esecuzione}\label{app:istruzioni}
% ═══════════════════════════════════════════════════════════════════════════

In questa appendice sono raccolte le istruzioni per eseguire il codice
Python presentato nell'Appendice~\ref{app:listati} e per avviare
l'applicazione web interattiva descritta nella Sezione~\ref{sec:architettura}.

\subsection{Dipendenze}

Il codice richiede Python 3.9 o successivo e le seguenti librerie:

\begin{itemize}
  \item \texttt{numpy} (versione $\geq 1.21$) -- calcolo numerico vettoriale;
  \item \texttt{matplotlib} (versione $\geq 3.5$) -- generazione dei grafici
  di convergenza e dei percorsi di ottimizzazione;
  \item \texttt{plotly} (versione $\geq 5.0$) -- visualizzazioni interattive
  nell'applicazione web;
  \item \texttt{pyodide} -- runtime Python nel browser per l'esecuzione del
  codice degli algoritmi (usato dall'applicazione web).
\end{itemize}

\subsection{Installazione}

Si consiglia di creare un ambiente virtuale dedicato e installare le
dipendenze con \texttt{pip}:

\begin{verbatim}
python3 -m venv .venv
source .venv/bin/activate        # su Windows: .venv\Scripts\activate
pip install --upgrade pip
pip install numpy matplotlib plotly
\end{verbatim}

\subsection{Esecuzione degli esperimenti numerici}

Gli esperimenti del Capitolo~\ref{sec:esperimenti} (convergenza, dinamica del
batch e confronto tra metodi) possono essere riprodotti eseguendo lo script
\texttt{sim\_exp.py}, che genera automaticamente le figure in formato PDF/PNG e
le tabelle dei risultati in formato LaTeX:

\begin{verbatim}
python3 sim_exp.py
\end{verbatim}

Lo script accetta i parametri descritti nel Capitolo~\ref{sec:esperimenti}
(numero di esempi $N$, dimensionalità $m$, tolleranza $\theta$, rapporto
$R$) e stampa a video, per ogni funzione test, il numero di iterazioni, il
tempo di esecuzione, l'errore finale $\|w_k - w^*\|$ e la dimensione massima
del batch raggiunta.

\subsection{Esecuzione dei singoli algoritmi}

I listati dell'Appendice~\ref{app:listati} possono essere eseguiti anche in
modo isolato, previa definizione delle funzioni ausiliarie
\texttt{loss\_i}, \texttt{grad\_i}, \texttt{hessvec\_i}, \texttt{grad\_full}
e della costante $N$, che dipendono dal dataset. Un esempio completo per la
regressione lineare è incluso nello script \texttt{sim\_exp.py}.

\subsection{Avvio dell'applicazione web}

L'applicazione web interattiva è un'applicazione a pagina singola (HTML +
JavaScript) che esegue il codice Python nel browser tramite Pyodide: non è
necessario un server dedicato. È sufficiente aprire il file HTML in un
browser moderno (Chrome, Firefox, Safari, Edge). Per lo sviluppo locale si può
utilizzare un semplice server HTTP statico:

\begin{verbatim}
python3 -m http.server 8000
\end{verbatim}

e quindi visitare \texttt{http://localhost:8000}.

\subsection{Requisiti di sistema}

Gli esperimenti riportati in questo elaborato sono stati eseguiti su un
computer portatile con processore Apple Silicon (8 core) e 16~GB di RAM, in
ambiente macOS, utilizzando Python 3.14, NumPy 2.4 e Matplotlib 3.11. Il costo
complessivo della simulazione (tre problemi test, tre valori di $\theta$ e
quattro algoritmi) è inferiore a due minuti; nessun acceleratore GPU è
richiesto, poiché tutti i calcoli sono operazioni vettoriali su array NumPy di
dimensione al più $N \times m$.
